\chapter{Fundamentos de IA e Modelos de Linguagem de Grande Escala}
\label{chap:fundamentos_llm_ia}

% ---------------------------------------------------------------
\section{Escopo do cap\'itulo}
% ---------------------------------------------------------------

Este cap\'itulo \'e o primeiro ato da hist\'oria que este livro conta.
Antes de qualquer agente, pipeline ou avalia\c{c}\~ao, \'e preciso
compreender a pe\c{c}a central do sistema: o modelo de linguagem que
recebe um \textit{prompt} e produz uma resposta. Tudo que vem depois
--- RAG, mem\'oria, ferramentas, juízes --- \'e uma forma de dar
melhores in\-for\-ma\-\c{c}\~oes a esse modelo ou de avaliar o que ele
produziu.

Este cap\'itulo apresenta os conceitos fundamentais que sustentam toda a
arquitetura desenvolvida ao longo do livro: modelos de linguagem de grande
escala (do ingl\^es \textit{Large Language Models} --- LLMs), mecanismos de
aten\c{c}\~ao, tokeniza\c{c}\~ao, janelas de contexto e engenharia de
\textit{prompts}. Ao dominar esses conceitos, o leitor adquire a base
te\'orica necess\'aria para compreender as decis\~oes de projeto adotadas
nos cap\'itulos subsequentes, em especial a escolha do modelo
\textit{Claude Sonnet} via \textit{Amazon Bedrock} como n\'ucleo de infer\^encia
da plataforma.

\begin{quote}\itshape
\textbf{Onde estamos na hist\'oria:} loop interno, fase de fundamentos.
O sistema ainda n\~ao existe --- estamos aprendendo a linguagem com a
qual ele ser\'a constru\'ido.
\end{quote}

% ---------------------------------------------------------------
\section{Objetivos de aprendizagem}
% ---------------------------------------------------------------

Ao concluir este cap\'itulo, o leitor ser\'a capaz de:

\begin{itemize}[leftmargin=*]
  \item Descrever a arquitetura \textit{Transformer} e o papel do mecanismo
        de aten\c{c}\~ao na gera\c{c}\~ao de texto;
  \item Explicar o processo de tokeniza\c{c}\~ao e calcular o consumo
        estimado de \textit{tokens} de uma requisi\c{c}\~ao;
  \item Formular \textit{prompts} eficazes utilizando as t\'ecnicas
        \textit{zero-shot}, \textit{few-shot} e
        \textit{chain-of-thought} (CoT)\footnote{%
          \textit{Chain-of-Thought} \'e uma t\'ecnica de
          \textit{prompting} que instrui o modelo a exibir o racioc\'inio
          passo a passo antes de produzir a resposta final
          \cite{wei2022chain}.};
  \item Identificar os par\^ametros de amostragem --- temperatura,
        \textit{top-p} e \textit{top-k} --- e seu impacto na qualidade
        e na diversidade das respostas;
  \item Invocar o modelo \textit{Claude Sonnet} via API
        \textit{Amazon Bedrock} a partir de c\'odigo Python.
\end{itemize}

% ---------------------------------------------------------------
\section{Conceitos te\'oricos}
% ---------------------------------------------------------------

\subsection{Arquitetura Transformer}

A arquitetura \textit{Transformer}\index{Transformer|textbf}\index{arquitetura!Transformer},
proposta por \citeonline{vaswani2017attention},
revolucionou o processamento de linguagem natural\index{PLN|see{processamento de linguagem natural}}\index{processamento de linguagem natural}
ao substituir redes recorrentes\index{redes recorrentes}
pelo mecanismo de \textit{self-attention}\index{self-attention|textbf}\index{aten\c{c}\~ao!self-attention},
que permite ao modelo relacionar
quaisquer dois tokens da sequ\^encia de entrada independentemente de sua
dist\^ancia posicional. Formalmente, a aten\c{c}\~ao \'e calculada como:

\begin{equation}
  \text{Attention}(Q, K, V) = \text{softmax}\!\left(\frac{QK^\top}{\sqrt{d_k}}\right)V
  \label{eq:attention}
\end{equation}

\noindent em que $Q$, $K$ e $V$ s\~ao as matrizes de consulta
(\textit{query}), chave (\textit{key}) e valor (\textit{value}),
respectivamente, e $d_k$ \'e a dimens\~ao da chave, utilizada como fator
de normaliza\c{c}\~ao para estabilizar os gradientes durante o treinamento.
A fun\c{c}\~ao \textit{softmax}\index{softmax|textbf}\footnote{%
  $\sigma(\mathbf{z})_i = e^{z_i}/\!\sum_j e^{z_j}$ --- converte o vetor
  de \textit{logits} em distribui\c{c}\~ao de probabilidade. Defini\c{c}\~ao
  formal no Ap\^endice~\ref{apendice:glossario_matematico},
  p.~\pageref{sec:softmax}.} normaliza os produtos escalados, produzindo
pesos de aten\c{c}\~ao que somam 1. O empilhamento de m\'ultiplas camadas de
aten\c{c}\~ao multi-cabe\c{c}a\index{aten\c{c}\~ao!multi-head}
(\textit{multi-head attention}) confere ao modelo a capacidade de capturar
padr\~oes sint\'aticos, sem\^anticos e pragm\'aticos simultaneamente.

LLMs s\~ao modelos \textit{Transformer} de \'unico-encoder-decoder ou
somente-decoder\index{decoder} treinados em corpora massivos mediante o objetivo de
previs\~ao do pr\'oximo \textit{token}\index{token} (\textit{next-token prediction}\index{next-token prediction}).
O volume de dados e par\^ametros --- tipicamente da ordem de bilh\~oes
--- confere a esses modelos capacidades emergentes n\~ao observadas em
escalas menores, como racioc\'inio l\'ogico, tradu\c{c}\~ao sem
supervis\~ao expl\'icita e gera\c{c}\~ao de c\'odigo
\cite{wei2022emergent}.

\subsection{Tokeniza\c{c}\~ao}

\textit{Tokens}\index{token} s\~ao as unidades m\'inimas de processamento de um LLM,
obtidos por meio de algoritmos de segmenta\c{c}\~ao subpalavra\index{segmenta\c{c}\~ao subpalavra}, como
\textit{Byte-Pair Encoding}\index{Byte-Pair Encoding|textbf}\index{BPE|see{Byte-Pair Encoding}} (BPE)\footnote{%
  Defini\c{c}\~ao formal e pseudoc\'odigo no
  Ap\^endice~\ref{apendice:glossario_matematico}, p.~\pageref{sec:bpe}.} e
\textit{SentencePiece}\index{SentencePiece}\footnote{%
  Defini\c{c}\~ao e compara\c{c}\~ao com BPE no
  Ap\^endice~\ref{apendice:glossario_matematico}, p.~\pageref{sec:sentencepiece}.}. A rela\c{c}\~ao
entre palavras e \textit{tokens} varia conforme o vocabul\'ario do modelo;
como regra pr\'atica, estima-se que 1.000 palavras em portugu\^es equivalem
a aproximadamente 1.300--1.500 \textit{tokens} no modelo \textit{Claude
Sonnet}.

O custo operacional e a lat\^encia de infer\^encia s\~ao diretamente
proporcionais ao n\'umero de \textit{tokens} processados, tornando a
contagem pr\'evia um fator cr\'itico no dimensionamento da plataforma.
No projeto descrito neste livro, adota-se a biblioteca \texttt{tiktoken}
\cite{tiktoken2023} como estimador de \textit{tokens} nos benchmarks de
desempenho apresentados no Cap\'itulo~\ref{chap:avaliacao_metricas_mercado}.

\subsection{Janela de Contexto}

A janela de contexto\index{janela de contexto|textbf} (\textit{context window}\index{context window|see{janela de contexto}}) define o n\'umero m\'aximo
de \textit{tokens} --- somando-se entrada e sa\'ida --- que o modelo
processa em uma \'unica requisi\c{c}\~ao. O \textit{Claude Sonnet~4.5},
utilizado na plataforma, suporta at\'e 200.000 \textit{tokens} de entrada,
o que possibilita a ingest\~ao de documentos extensos sem necessidade de
segmenta\c{c}\~ao agressiva. Todavia, o custo por \textit{token} de entrada
\'e inferior ao de sa\'ida; portanto, o projeto limita a janela de
contexto injetada pelo pipeline RAG --- descrito no
Cap\'itulo~\ref{chap:rag_fundamentos} --- a 8.000 \textit{tokens}, valor
determinado empiricamente como suficiente para manter alta qualidade de
resposta sem comprometer a lat\^encia.

\subsection{Par\^ametros de Amostragem}

A distribui\c{c}\~ao de probabilidade sobre o vocabul\'ario na etapa de
decodifica\c{c}\~ao\index{decodifica\c{c}\~ao} \'e modulada por tr\^es par\^ametros principais:

\begin{description}
  \item[Temperatura ($\tau$)\index{temperatura (LLM)|textbf}] Reescala os logits\index{logits}\footnote{%
        \textit{Logits}: valores reais n\~ao normalizados produzidos pela
        \'ultima camada linear do modelo. Defini\c{c}\~ao formal no
        Ap\^endice~\ref{apendice:glossario_matematico}, p.~\pageref{sec:logits}.}
        antes do
        \textit{softmax}\index{softmax}\footnote{%
        Fun\c{c}\~ao $\sigma(\mathbf{z})_i = e^{z_i}/\sum_j e^{z_j}$
        que converte logits em distribui\c{c}\~ao de probabilidade. Defini\c{c}\~ao
        completa no Ap\^endice~\ref{apendice:glossario_matematico},
        p.~\pageref{sec:softmax}.}: valores baixos ($\tau \to 0$) tornam a
        distribui\c{c}\~ao mais pontuda, favorecendo respostas
        deterministas; valores altos ($\tau > 1$) aumentam a entropia
        e, portanto, a diversidade da gera\c{c}\~ao.
  \item[\textit{Top-p}\index{top-p|textbf}\index{amostragem por n\'ucleo|see{top-p}} (amostragem por n\'ucleo)] Restringe o
        vocabul\'ario ativo ao menor conjunto de \textit{tokens} cuja
        probabilidade acumulada supera $p$; valores t\'ipicos situam-se
        entre $0,85$ e $0,95$.
  \item[\textit{Top-k}\index{top-k|textbf}] Limita a amostragem aos $k$ \textit{tokens}
        mais prov\'aveis; \'e frequentemente combinado com
        \textit{top-p} para eliminar candidatos de baixa probabilidade.
\end{description}

No projeto, adota-se $\tau = 0{,}3$ nas consultas RAG de modo
\texttt{segmento} e \texttt{persona}, priorizando consist\^encia
e fidelidade ao contexto recuperado. Consultas de modo
\texttt{twin} utilizam $\tau = 0{,}5$ para permitir maior
variabilidade na reprodu\c{c}\~ao do comportamento individual.

% ---------------------------------------------------------------
\section{Implementa\c{c}\~ao orientada por passos}
% ---------------------------------------------------------------

\subsection{Passo 1 --- Configura\c{c}\~ao das credenciais AWS}

O acesso ao \textit{Amazon Bedrock}\index{Amazon Bedrock|textbf}\index{Bedrock|see{Amazon Bedrock}} requer credenciais IAM\index{IAM}\index{Identity and Access Management|see{IAM}} com a permiss\~ao
\texttt{bedrock:InvokeModel}. Recomenda-se utilizar perfis nomeados via
\texttt{AWS CLI} para evitar a gest\~ao manual de chaves no c\'odigo:

\begin{lstlisting}[language=bash, caption={Configura\c{c}\~ao de perfil AWS CLI.}]
aws configure --profile langchain-dev
# AWS Access Key ID: <ID>
# AWS Secret Access Key: <CHAVE>
# Default region name: us-east-1
# Default output format: json
\end{lstlisting}

\subsection{Passo 2 --- Primeira chamada ao modelo}

O cliente Bedrock \'e instanciado via \texttt{boto3}, e a
requisi\c{c}\~ao \'e formatada conforme a API \textit{Messages} da
Anthropic:

\begin{lstlisting}[language=Python,
  caption={Primeira chamada ao \textit{Claude Sonnet} via Bedrock.},
  label={lst:bedrock_hello}]
import boto3, json

MODEL_ID = "us.anthropic.claude-sonnet-4-5-20250514-v1:0"
REGION   = "us-east-1"

client = boto3.client("bedrock-runtime", region_name=REGION)

payload = {
    "anthropic_version": "bedrock-2023-05-31",
    "max_tokens": 512,
    "temperature": 0.3,
    "messages": [
        {
            "role": "user",
            "content": "Explique o conceito de tokenizacao em LLMs "
                       "em ate tres paragrafos, estilo academico."
        }
    ]
}

response = client.invoke_model(
    modelId=MODEL_ID,
    body=json.dumps(payload),
    contentType="application/json",
    accept="application/json"
)

body   = json.loads(response["body"].read())
texto  = body["content"][0]["text"]
tokens = body["usage"]          # input_tokens, output_tokens

print(texto)
print(f"\nTokens consumidos: entrada={tokens['input_tokens']}, "
      f"saida={tokens['output_tokens']}")
\end{lstlisting}

Verifica-se que o campo \texttt{usage} retornado pela API cont\'em a
contagem exata de \textit{tokens} de entrada e sa\'ida, permitindo o
rastreamento preciso dos custos de infer\^encia.

\subsection{Passo 3 --- Engenharia de \textit{prompts}}

A qualidade da resposta \'e sens\'ivel \`a estrutura do \textit{prompt}.
O padr\~ao adotado no projeto separa o \textit{system prompt}\index{prompt!system prompt|textbf} --- que
define o papel, o tom e as restri\c{c}\~oes do modelo --- das mensagens
de usu\'ario, conforme recomendado pela documenta\c{c}\~ao da Anthropic
\cite{anthropic2024prompting}:

\begin{lstlisting}[language=Python,
  caption={Estrutura de \textit{system prompt} e mensagem de usu\'ario.},
  label={lst:prompt_structure}]
payload = {
    "anthropic_version": "bedrock-2023-05-31",
    "max_tokens": 1024,
    "temperature": 0.3,
    "system": (
        "Voce e um assistente especializado em analise de clientes "
        "bancarios. Responda em portugues do Brasil, estilo formal, "
        "com base exclusivamente no contexto fornecido. "
        "Nao invente informacoes ausentes do contexto."
    ),
    "messages": [
        {
            "role": "user",
            "content": (
                "<context>\n{contexto_recuperado}\n</context>\n\n"
                "Pergunta: {pergunta_usuario}"
            )
        }
    ]
}
\end{lstlisting}

O uso de delimitadores XML (\texttt{<context>}) \'e recomendado para
que o modelo distinga claramente o contexto recuperado da
instru\c{c}\~ao, reduzindo o risco de alucina\c{c}\~oes.

\subsection{Passo 4 --- T\'ecnica \textit{few-shot}}

Exemplos de entrada--sa\'ida inseridos no \textit{prompt}
(\textit{few-shot examples}\index{few-shot|textbf}\index{prompt!few-shot}) calibram o formato e o registro das
respostas sem necessidade de ajuste fino\index{fine-tuning|textbf} (\textit{fine-tuning}). No
projeto, cada modo de consulta --- \texttt{segmento}, \texttt{persona}
e \texttt{twin} --- utiliza de dois a tr\^es exemplos curados que
demonstram a estrutura esperada da resposta, conforme detalhado no
Cap\'itulo~\ref{chap:segmentacao_personas_twins}.

% ---------------------------------------------------------------
\section{Plano de testes do cap\'itulo}
% ---------------------------------------------------------------

\subsection{Teste funcional}

\begin{itemize}[leftmargin=*]
  \item \textbf{TC-01:} Invocar o modelo com \textit{prompt} m\'inimo
        e verificar que o campo \texttt{content[0].text} est\'a
        preenchido e n\~ao \'e vazio.
  \item \textbf{TC-02:} Passar um \textit{prompt} com temperatura
        $\tau = 0$ e verificar que duas execu\c{c}\~oes consecutivas
        produzem sa\'idas id\^enticas (determinismo).
  \item \textbf{TC-03:} Exceder o limite de \textit{max\_tokens} e
        confirmar que a resposta \'e truncada com o motivo
        \texttt{max\_tokens} no campo \texttt{stop\_reason}.
\end{itemize}

\subsection{Teste de desempenho}

\begin{itemize}[leftmargin=*]
  \item \textbf{TP-01:} Medir a lat\^encia de 10 chamadas
        consecutivas com \textit{prompt} de 500 \textit{tokens};
        crit\'erio: mediana $\leq 3{,}0$\,s (regi\~ao
        \texttt{us-east-1}).
  \item \textbf{TP-02:} Calcular o custo estimado por 1.000
        requisi\c{c}\~oes de 1.000 \textit{tokens} de entrada e
        500 de sa\'ida com base na tabela de pre\c{c}os
        \textit{on-demand} do Bedrock.
\end{itemize}

\subsection{Teste de qualidade de resposta}

\begin{itemize}[leftmargin=*]
  \item \textbf{TQ-01:} Formular cinco perguntas de refer\^encia
        sobre tokeniza\c{c}\~ao e avaliar as respostas quanto \`a
        precis\~ao factual (escala 1--5).
  \item \textbf{TQ-02:} Verificar que a resposta a uma pergunta
        fora do contexto retorna a express\~ao de aus\^encia de
        informa\c{c}\~ao, sem alucina\c{c}\~ao.
\end{itemize}

% ---------------------------------------------------------------
\section{Crit\'erios de aceite}
% ---------------------------------------------------------------

\begin{itemize}[leftmargin=*]
  \item TC-01, TC-02 e TC-03 executam sem erros bloqueantes.
  \item TP-01: mediana de lat\^encia $\leq 3{,}0$\,s confirmada em
        execu\c{c}\~ao real (evidenciar com \textit{log} de timestamps).
  \item TQ-01: pontua\c{c}\~ao m\'edia $\geq 4{,}0/5{,}0$.
  \item TQ-02: zero alucina\c{c}\~oes detectadas nas cinco perguntas
        fora de contexto.
\end{itemize}

% ---------------------------------------------------------------
\section{Espa\c{c}o para expans\~ao iterativa}
% ---------------------------------------------------------------

\begin{itemize}[leftmargin=*]
  \item \textbf{v1.1} --- Adicionar compara\c{c}\~ao de modelos:
        \textit{Claude Haiku} (baixo custo) vs. \textit{Claude Sonnet}
        (equil\'ibrio) vs. \textit{Claude Opus} (m\'axima qualidade),
        com m\'etricas de lat\^encia e custo por token.
  \item \textbf{v1.2} --- Incluir diagrama da arquitetura
        \textit{Transformer} (codificador e decodificador) adaptado
        ao estilo do livro.
  \item \textbf{Lacuna} --- Avaliar impacto de diferentes estrat\'egias
        de truncamento de contexto (\textit{sliding window} vs.
        sumariza\c{c}\~ao hier\'arquica) na qualidade das respostas.
\end{itemize}
